\section{Discussion of Findings: Insights into Current Smart Building Research}
\label{sec-ch2:implications}
We now discuss our findings to elicit higher-level understandings and implications. The aim is to extract messages that can inform the growing branch of research that seeks to design human experiences in smart buildings.

\subsection{Beyond Comfort---The Diversity of Targeted Human Experiences in Smart Building Research}
\label{sec-ch2:implications-landscape}
In our exploration of human experiences in smart buildings, we identified 11 distinct experiences around human experiences in smart building research. This reflects an evolving understanding of what it means to design a smart building and extends it substantially beyond the traditionally examined experience of comfort. For example, \textbf{Learning and Teaching Experiences} highlight the potential of smart buildings as educational tools that create interactive learning environments which support curiosity and innovation. \cite{lui2014supporting}. Although \textbf{Comfort} remains a major focus, it is often addressed through novel means such as \textit{Personalisation}~\cite{radic2023advanced} and  \textit{Spatial Design}~\cite{nguyen2022towards}, as well as fostering \textit{Community and Social Connections}~\cite{olsen2017wizard} to target individual and collective well-being. 
This diversity of targeted human experiences beyond comfort sheds light on the major challenges and questions that need to be addressed for the development of the field. For example, what design methods must we develop or adopt to conceptualise emotionally-rich~\cite{rao2023towards} or ``somatic'' environments ~\cite{alfonso2024embodied}? How can we design for meaningful yet highly subjective \textbf{Inclusion Experiences} in smart buildings, catering to those with diverse sensory needs~\cite{swaminathan2021tactile, ahlquist2017sensory}? The same question is valid for evaluation methods to assess the impact of technological interventions on these targeted human experiences. 


\subsection{Alignment with Trends in HCI}
Our analysis reveals strong connections between human experience design in smart buildings and long-established HCI research, particularly in the use of \textit{Community and Social Connection} as a key design measure for the \textbf{Experience of Relatedness} among occupants. Smart buildings can foster social connection between occupants through technologies like multimodal communication systems for remote users~\cite{hansen2009pendaphonics} and features that enhance co-located social bonding~\cite{nabil2020peace}. These works emphasise the potential of smart buildings to bridge physical and digital spaces, enhancing social connectedness across diverse contexts. However, our scoping review reveals that several targeted experiences in smart buildings also align with some of the more recent and growing trends in HCI research and related fields, including \textbf{Bodily Experiences}, \textbf{Explorative and/or Future and Speculative}, and \textbf{Democratic Experiences}. 

\textbf{Bodily Experiences} emphasise the role of multi-sensorial perception of the environment in shaping interactions within smart buildings, as well as mutual influences of the human body and spaces  that can be studied and designed for, using methods of \textit{Somaesthetic Design}. This builds on an architectural legacy in which the human body---its scale, directionality, proportionality, etc.---is the main source of inspiration and reference for spatial design \cite{shusterman2011somaesthetics}. We observe an alignment between the proliferation of works recognising the importance of bodily experiences in smart buildings with recent discourses and design instances surrounding Somaesthetic Design in HCI and IxD research. 

\textbf{Explorative and/or Future and Speculative Experiences} captures a body of work that aims to (re-)define how people experience and re-experience home and public spaces in novel and previously un-imagined ways by deploying design probes~\cite{gaver2004cultural} or speculative inquiry ~\cite{galloway2018speculative}. This experience is particularly relevant as it expands the concept of interaction beyond traditional interfaces, incorporating the \textit{Spatial Design} of the smart building itself as a medium for digital engagement. It aligns with the evolving discourse on architectonic interaction, which involves ``Interaction Design at the scale of architecture'' and the ``design of architectonic elements with interactive capabilities'' ~\cite{wiberg2020interaction}. As we increasingly occupy spaces where digital technologies work alongside us, this concept becomes essential for exploring how these environments are created and experienced.

Finally, \textbf{Democratic Experiences} reflects recent efforts in HCI to cultivate and preserve democratic values within the context of smart buildings \cite{dalton2020habitech}. This work views smart buildings as civic spaces that foster dialogue between occupants and their environment, enabling greater user agency, participation, and control, while empowering users to actively shape their surroundings \cite{dalton2020habitech}. This theme builds on well-established research in \textit{Participation and co-design}, while situating itself within the relatively new domains of smart cities and civic technologies, and critical design as mechanisms in HCI for co-shaping interactive spaces through occupant involvement.

\subsection{The Influences of Interaction Design and Architecture in Shaping Smart Building Research}
Our findings suggest that smart building research has been shaped by both IxD and architecture, driving the field forward in three ways: (1) by re-appropriating concepts such as personalisation that draw heavily from the traditions of IxD \cite{lundgaard2019temporal}, (2) by learning from space design principles rooted in the legacy of architecture~\cite{hillier1976space}, with a focus on the ability to ``experiment'' with how technology can introduce new forms of interactive experiences to our living spaces (e.g. through \textbf{Explorative and/or Future and Speculative} approaches \cite{takeuchi2009bezier}), and (3) by bringing together the two disciplines: while architecture provides the physical and cultural context, IxD enhances these spaces with interactive elements that foster connection and identity, e.g. enabling a \textbf{Sense of Placemaking}~\cite{alfonso2024embodied}. This fluidity between the two disciplines highlights the multidisciplinary nature of the smart building field which benefits from a collaborative approach to synthesising insights from multiple domains. By the same token, these observations raise important questions about what other knowledge could be transferred to smart building research. What concepts from other disciplines have yet to make their way into smart buildings? By exploring and opening new channels for interdisciplinary exchange, smart building research can continue to evolve, drawing on a broader range of ideas to shape the future of our smart buildings. For example, notions around digital rights (developed under the domain of digital philosophies~\cite{leonelli2016philosophy}) in online spaces could offer opportunities to address the challenges of smart buildings, which may collect data both voluntarily and involuntarily from occupants. Specifically, data portability~\cite{leonelli2016philosophy} could allow occupants to transfer their data preferences (e.g. lighting or temperature settings) between smart buildings, ensuring continuity and minimal data collection (as data gets collected once and re-used across different buildings~\cite{turner2021exercisability}). Similarly, smart buildings could enable ``guest modes'' (similar to online browsing experiences) to minimise data collection, privacy zones where sensors are deactivated, or informed consent mechanisms alongside transparency dashboards to display building data usage policies~\cite{santucci2013privacy}.



\subsection{The Integration of Artificial Intelligence in Human-Centred Smart Buildings}
Despite the growing adoption of artificial intelligence (AI) and its emergent applications in smart building design---such as building resource management, building information modelling \cite{coupry2021bim}, and building security \cite{genkin2023b}---there is very limited to almost no integration of technologies like generative AI and large language models (LLMs) for designing human experiences in smart buildings. When AI is integrated with human interaction, it primarily focuses on creating responsive and personalised environments. The use of AI was noted in \textbf{Comfort} experiences~\cite{yang2014making, zhao2017mediated} and considered in \textbf{Emotional Experiences}~\cite{rao2023towards}. Among the reviewed works, 14 utilised machine learning (ML) for example in tasks like adjusting temperature settings based on user behaviour \cite{yang2014making} (comfort). Other examples where ML is used include \citet{zhao2017mediated}'s \textit{``Mediated Atmospheres''} or \citet{andres2023human}'s \textit{``System of a Sound''} transform environments by digitally controlling lighting, sound, and video projections, creating spaces that interact with human emotions.

In addition, some works consider (but do not implement) the role of AI in futuristic concepts (like empathic buildings~\cite{rao2023towards}) that could adjust spaces to react to humans. For example, this extends to ideas like ``curious agents'' and ``motivated agents'' \cite{merrick2007curious}, which might support proactive problem-solving and adaptability within environments like ``resilient hospitals''  \cite{taylor2022hospitals}.


\subsection{``Reactive'' Retrofitted Design in Smart Buildings}
Our analysis around the technological interventions in smart buildings (\autoref{sec-ch2:findings-technology}) shows that the majority are concentrated in the \textit{Space Plan} and \textit{Stuff} layers \cite{brand1995buildings}, which are inherently more flexible and adaptable. The current distribution of technological interventions across the building layers thus reveals a reactive approach to interaction design in smart buildings. This late involvement raises questions about \emph{how interaction designers can get involved earlier in the process of planning and design of future buildings?} Early involvement of designers offers the potential to integrate interactive systems directly into the building’s fundamental layers, such as \textit{Structure} and \textit{Skin}, rather than retrofitting technology into more transient layers like \textit{Space Plan} and \textit{Stuff}. Designers, when included early, can employ iterative and participatory methods to ensure that the resulting smart buildings address the needs of occupants \cite{alavi2018hide}. However, challenges remain in achieving the early involvement of designers. As observed by \citet{verma2017studying}, the building design process is bound by rigid schedules and defined objectives, leaving little room for experimentation. According to \citet{verma2017studying}, the tendency to finalise most design decisions before construction begins makes the integration of interaction design particularly challenging. Furthermore, the ``siloed'' nature of architecture, engineering, and technology disciplines complicates cross-disciplinary collaboration where construction timelines and budget constraints often favour established workflows, creating resistance to adopting iterative and participatory strategies~\cite{verma2017studying}.
While mindful of these challenges, there is potential in integrating smart building technology into the building structure from the outset, rather than being retrofitted into the more transient layers.